\begin{center}
\Huge
Beviser
\end{center}
\section*{Sætning 1}
\stepcounter{section}

Vi skal vise, at der for funktionerne $f$ og $g$ givet ved
\begin{align*}
	f(x) = e^{kx}
\end{align*}
og 
\begin{align*}
	g(x) = a^x
\end{align*}
gælder, at $f$ og $g$ er differentiable funktioner, og at deres afledede funktioner er givet ved henholdsvis
\begin{align*}
	f'(x) = ke^{kx}
\end{align*}
og 
\begin{align*}
	g'(x) = \ln(a)a^x.
\end{align*}
\subsubsection*{Vejledning til $f'$:}
\begin{enumerate}[label = \roman*)]
	\item Afgør, hvad der tilsvarer den ydre og indre funktion for $f$
	\item Differentiér $f$. 
	\item Sæt firkanten.
\end{enumerate}

\subsubsection*{Vejledning til $g'$:}
\begin{enumerate}[label=\roman*)]
	\item Overbevis jer selv om, at $a = e^{\ln(a)}$ og indsæt dette på $a$'s plads i forskriften. 
	\item Anvend en potensregneregel til at omskrive udtrykket.
	\item Differentiér nu $g$.
	\item Sæt firkanten.
\end{enumerate}

\newpage
\section*{Sætning 2}
\stepcounter{section}

Vi skal vise, at funktionen $f:\mathbb{R}^+ \to \mathbb{R}$ givet ved
\begin{align*}
	f(x) = \ln(x)
\end{align*}
er differentiabel og har differentialkvotienten
\begin{align*}
	f'(x) = \frac{1}{x}
\end{align*}
\subsubsection*{Vejledning:}
\begin{enumerate}[label=\roman*)]
	\item Overbevis jer selv om, at 
	\begin{align}\label{eq:1}
		x = e^{\ln(x)}.
	\end{align}
	\item Differentiér begge sider af lighedstegnet i \eqref{eq:1}. Husk, at du ikke kan differentiere 
	$\ln(x)$,
	da vi skal bevise differentialkvotienten for netop denne. 
	\item Isolér $(\ln(x))'$ i udtrykket. 
	\item Sæt firkanten. 
\end{enumerate}

\newpage
\section*{Sætning 3}
\stepcounter{section}

Vi skal vise, at funktionen $f$ givet ved
\begin{align*}
	f(x) = x^a
\end{align*}
er differentiabel og har differentialkvotienten
\begin{align*}
	f'(x) = ax^{a-1}.
\end{align*}
\subsubsection*{Vejledning:}
\begin{enumerate}[label=\roman*)]
	\item Indsæt $x = e^{\ln(x)}$ på $x$'s plads i $f$. 
	\item Anvend en potensregnereglen på udtrykket for $f$.
	\item Differentiér udtrykket for $f$. 
	\item Brug tricket fra i) baglæns.
	\item Reducer nu på udtrykket.
	\item Sæt firkanten. 
\end{enumerate}

\newpage
\section*{Sætning 4}
\stepcounter{section}

Vi skal vise, at der for to differentiable funktioner $f$ og $g$, hvor $g(x) \neq 0$ gælder, at 
\begin{align*}
	 \left(\frac{f(x)}{g(x)} \right)' = \frac{f'(x)g(x)-f(x)g'(x)}{(g(x))^2}
\end{align*}
\subsubsection*{Vejledning:}
\begin{enumerate}[label=\roman*)]
	\item Opskriv
	\begin{align*}
		\frac{f(x)}{g(x)} = f(x)\cdot \frac{1}{g(x)}.
	\end{align*}
	 \item Brug kædereglen til at bestemme 
	 \begin{align*}
	 	\left(\frac{1}{g(x)}\right)'
	 \end{align*}
	 \item Brug nu produktreglen til at bestemme 
	 \begin{align*}
	 	\left( f(x) \cdot \frac{1}{g(x)}\right)'
	 \end{align*}
	 \item Forlæng brøken 
	 \begin{align*}
	 	f'(x)\cdot \frac{1}{g(x)}
	 \end{align*}
	 med $g(x)$ og sæt de to brøker på fælles brøkstreg.
	 \item Sæt firkanten. 
\end{enumerate}

\newpage
\section*{Sætning 5}
\stepcounter{section}

Vi skal vise, at vi kan skrive $\ln(a)$ som grænseværdien
\begin{align*}
	\ln(a) = \lim_{h \to 0}\frac{a^h-1}{h}.
\end{align*}
\subsubsection*{Vejledning:}
\begin{enumerate}[label=\roman*)]
	\item Opskriv sekanthældningen for $a^x$. (Brug $x$ i stedet for $x_0$.)
	\item Anvend en potensregneregel og sæt $a^x$ uden for parentes. 
	\item Tag grænseværdien $h \to 0$. 
	\item Sæt $a^x$ uden for grænseværdien. 
	\item Udnyt, at $(a^x)' = \ln(a)a^x$, som du viste i Sætning 1. 
	\item Lad $(a^x)'$ gå ud og sæt firkanten. 
\end{enumerate}